\chapter{Mental Technique}\label{ch:iv1:mental-technique}

\begin{ivtrains}
The methods behind fast arithmetic: why a shortcut works, when it fails and how to check its answer. Interviewers ask about these directly, and they also expect them along the way in longer questions. The chapter starts with products and percentages, moves to compounding, logs and roots, and ends with estimates that combine them. No calculator, and say each step aloud.
\end{ivtrains}

\begin{question}[$\star$ \iqroles{QT, Trader, Analyst} \iqfirm{any} \iqlevel{junior}]\label{iq:iv1:mental-technique:1}
In your head, what is 97 times 103? What makes it quick, and when does the same idea work?
\end{question}

\begin{question}[$\star$ \iqroles{QT, Trader} \iqfirm{market maker, bank} \iqlevel{junior}]\label{iq:iv1:mental-technique:2}
What is 85 squared? There is a quick rule for squaring any number that ends in 5: state it and prove it.
\end{question}

\begin{question}[$\star$ \iqroles{QT, Trader} \iqfirm{market maker} \iqlevel{junior}]\label{iq:iv1:mental-technique:3}
Both 97 and 96 are just below 100. Show how to multiply them using only how far each one is from 100, and explain why your method works.
\end{question}

\begin{question}[$\star$ \iqroles{QT, Trader, Analyst} \iqfirm{any} \iqlevel{junior}]\label{iq:iv1:mental-technique:4}
What is 368 times 125, in your head? What general trick are you using, and which other multipliers does it cover?
\end{question}

\begin{question}[$\star$ \iqroles{QT, Trader} \iqfirm{market maker, bank} \iqlevel{junior}]\label{iq:iv1:mental-technique:5}
There is a trick for multiplying by 11 by adding neighbouring digits. Use it on 486 times 11, including the carries, and explain why it works.
\end{question}

\begin{question}[$\star$ \iqroles{QT, Trader, Analyst, Risk} \iqfirm{any} \iqlevel{junior}]\label{iq:iv1:mental-technique:6}
What is 16\% of 25? Why can you swap the two numbers in a question like this? Use the same idea on 18\% of 50, and on 4\% of 175.
\end{question}

\begin{question}[$\star$ \iqroles{QT, Trader, Analyst, Risk} \iqfirm{any} \iqlevel{junior}]\label{iq:iv1:mental-technique:7}
A price of 384 includes 20\% tax on the pre-tax price. How much of the 384 is tax? Give the quickest method.
\end{question}

\begin{question}[$\star$ \iqroles{QT, Trader, Analyst, Risk} \iqfirm{any} \iqlevel{junior}]\label{iq:iv1:mental-technique:8}
An investment falls 20\%. What percentage rise brings it back to where it started? Give the general formula for a fall of any size, and use it for a fall of 37.5\%.
\end{question}

\begin{question}[$\star$ \iqroles{QT, Trader, Analyst} \iqfirm{any} \iqlevel{junior}]\label{iq:iv1:mental-technique:9}
Which is larger, $7/9$ or $11/14$, and by how much? Answer without turning either into a decimal.
\end{question}

\begin{question}[$\star$ \iqroles{QT, Trader} \iqfirm{market maker} \iqlevel{junior}]\label{iq:iv1:mental-technique:10}
You know that $1/7 = 0.142857142857\dots$, with those six digits repeating. How do you write $3/7$ and $5/7$ as decimals without dividing, and why does your method work?
\end{question}

\begin{question}[$\star$ \iqroles{QT, Trader, Analyst} \iqfirm{any} \iqlevel{junior}]\label{iq:iv1:mental-technique:11}
You buy 300 shares at 80 and 500 shares at 88. What is your average price? Find it without working out either total cost, and say what rule you are using.
\end{question}

\begin{question}[$\star$ \iqroles{QT, Trader, Analyst, Risk} \iqfirm{any} \iqlevel{junior}]\label{iq:iv1:mental-technique:12}
By the rule of 72, how long does money take to double at 8\% a year, compounded yearly? How close is that to the exact answer, and why does a rule like this work at all?
\end{question}

\begin{question}[$\star$ \iqroles{QT, Trader, Risk} \iqfirm{any} \iqlevel{junior}]\label{iq:iv1:mental-technique:13}
A stock has an annual volatility of 32\%, meaning that its yearly return has a standard deviation of 32\%. There are 252 trading days in a year, and daily returns are independent. What is the standard deviation of one day's return? What quick number do you divide by?
\end{question}

\begin{question}[$\star$ \iqroles{QR, QT, Trader} \iqfirm{any} \iqlevel{junior}]\label{iq:iv1:mental-technique:14}
You are told only that $\log_{10} 2 \approx 0.301$ and $\log_{10} 3 \approx 0.477$. Find $\log_{10} 72$ and $\log_{10} 0.0072$ without a calculator.
\end{question}

\begin{question}[$\star$ \iqroles{QT, QD, Dev, LLE} \iqfirm{any} \iqlevel{junior}]\label{iq:iv1:mental-technique:15}
Without a calculator, roughly how large is $2^{30}$, and how many doublings take you from 1 to one trillion? How far off is your first estimate, and what fact are you using?
\end{question}

\begin{question}[$\star$ \iqroles{QR, QT, Trader} \iqfirm{any} \iqlevel{junior}]\label{iq:iv1:mental-technique:16}
Estimate $\sqrt{50}$ to three decimal places with one step of mental work. Is your estimate too high or too low, and why is it always on that side?
\end{question}

\begin{question}[$\star$ \iqroles{QT, Trader, Analyst, Risk} \iqfirm{bank, hedge fund} \iqlevel{junior}]\label{iq:iv1:mental-technique:17}
One basis point is one hundredth of one percent. What is 10 basis points of \$3.5 million? A loan of \$250 million pays an extra one basis point a year: what is that in dollars over three months? Give the quick method.
\end{question}

\begin{question}[$\star$ \iqroles{QT, Trader, Risk} \iqfirm{bank, hedge fund} \iqlevel{junior}]\label{iq:iv1:mental-technique:18}
A bond price is quoted in 32nds of a point, so that 99-17 means 99 and $17/32$. What is 99-17 as a decimal, and how do you get it without dividing by 32?
\end{question}

\begin{question}[$\star\star$ \iqroles{QR, QT, Trader} \iqfirm{any} \iqlevel{junior}]\label{iq:iv1:mental-technique:19}
Without working out either product in full, which is larger, $37 \times 43$ or $38 \times 42$, and by exactly how much? What is the general rule?
\end{question}

\begin{question}[$\star\star$ \iqroles{QT, Trader, Risk} \iqfirm{any} \iqlevel{junior}]\label{iq:iv1:mental-technique:20}
Work out $1.02 \times 0.97$ to four decimal places, in your head. Which part of the calculation may you drop when you only need a rough answer, and how large is the error?
\end{question}

\begin{question}[$\star\star$ \iqroles{QR, QT, Trader, QD, Dev} \iqfirm{any} \iqlevel{junior}]\label{iq:iv1:mental-technique:21}
Someone tells you that $3847 \times 26 = 100\,122$. Show that this is wrong without doing the multiplication. Which quick checks does it pass, and which does it fail?
\end{question}

\begin{question}[$\star\star$ \iqroles{QT, Trader} \iqfirm{market maker} \iqlevel{junior}]\label{iq:iv1:mental-technique:22}
What is $73 \times 68$: 4964, 4954, 5064 or 4864? All four end in the right digit and have about the right size. How do you pick the answer without multiplying?
\end{question}

\begin{question}[$\star\star$ \iqroles{QR, QT, QD} \iqfirm{any} \iqlevel{junior}]\label{iq:iv1:mental-technique:23}
Someone writes $71 \times 88 = 6428$. Adding up the digits and comparing modulo 9 does not catch this error. Why not, and what quick check does catch it?
\end{question}

\begin{question}[$\star\star$ \iqroles{QT, Trader, Analyst, Risk} \iqfirm{any} \iqlevel{junior}]\label{iq:iv1:mental-technique:24}
A price rises 20\% and then falls 15\%. What is the net change, and does the order matter? Give a rule for combining two percentage moves in your head.
\end{question}

\begin{question}[$\star\star$ \iqroles{QR, QT, Trader} \iqfirm{any} \iqlevel{junior}]\label{iq:iv1:mental-technique:25}
Estimate $1/0.97$ and $1/1.03$ to four decimal places without dividing. What approximation do you use, and when does it stop working?
\end{question}

\begin{question}[$\star\star$ \iqroles{QR, QT, Trader, Risk} \iqfirm{any} \iqlevel{junior}]\label{iq:iv1:mental-technique:26}
Estimate $1.05^{10}$ to two decimal places, in your head, in two different ways.
\end{question}

\begin{question}[$\star\star$ \iqroles{QR, QT, Trader, Risk} \iqfirm{market maker, hedge fund} \iqlevel{junior}]\label{iq:iv1:mental-technique:27}
Estimate $0.99^{100}$ to three decimal places without a calculator. You may use $1/e \approx 0.368$. Why is $1/e$ alone not good enough?
\end{question}

\begin{question}[$\star\star$ \iqroles{QR, QT, Trader} \iqfirm{any} \iqlevel{junior}]\label{iq:iv1:mental-technique:28}
The rule of 72 says that money doubles in about $72/r$ years at $r$ percent a year, compounded yearly. How wrong is it at $r = 1$ and at $r = 36$? Near which rate is it exact? What simple correction fixes it at all these rates?
\end{question}

\begin{question}[$\star\star$ \iqroles{QR, QT, Trader, Risk} \iqfirm{market maker, hedge fund} \iqlevel{junior}]\label{iq:iv1:mental-technique:29}
Your account grows by 10 basis points every day for 252 days, where one basis point is 0.01\%. Estimate the total return in your head to one decimal place of a percent. What is the simple answer, and why is it too low?
\end{question}

\begin{question}[$\star\star$ \iqroles{QR, QT, Trader, Risk} \iqfirm{market maker, hedge fund} \iqlevel{junior}]\label{iq:iv1:mental-technique:30}
Your investment rises 10\% and falls 10\% in turn, for 100 periods. What fraction of your money is left at the end? The average period's return is zero, so why have you lost money? State the rule behind the answer.
\end{question}

\begin{question}[$\star\star$ \iqroles{QT, Trader, Risk} \iqfirm{any} \iqlevel{junior}]\label{iq:iv1:mental-technique:31}
A position loses 3\% of its remaining value each year. Roughly how long until half is gone? Does the rule of 72 give too long or too short a time here, and how does this differ from the doubling case at the same rate?
\end{question}

\begin{question}[$\star\star$ \iqroles{QR, QT} \iqfirm{any} \iqlevel{junior}]\label{iq:iv1:mental-technique:32}
You are told only that $\log_{10} 2 \approx 0.301$ and $\log_{10} 3 \approx 0.477$. Find $\log_{10} 7$ to three decimal places without a calculator.
\end{question}

\begin{question}[$\star\star$ \iqroles{QR, QT, QD, Dev} \iqfirm{any} \iqlevel{junior}]\label{iq:iv1:mental-technique:33}
How many decimal digits does $2^{100}$ have, and what is its first digit? You may use $\log_{10} 2 \approx 0.30103$.
\end{question}

\begin{question}[$\star\star$ \iqroles{QR, QT, Trader} \iqfirm{any} \iqlevel{junior}]\label{iq:iv1:mental-technique:34}
You know that $\ln 2 \approx 0.693$ and $\ln 10 \approx 2.303$. Estimate $\ln 1.3$ to three decimal places. Why does the series for $\ln(1 + x)$ alone serve you badly here?
\end{question}

\begin{question}[$\star\star$ \iqroles{QR, QT, Trader} \iqfirm{any} \iqlevel{junior}]\label{iq:iv1:mental-technique:35}
Estimate the tenth root of 2 and the fifth root of 3, each to three decimal places, without a calculator. You may use $\ln 2 \approx 0.693$ and $\ln 3 \approx 1.099$.
\end{question}

\begin{question}[$\star\star$ \iqroles{QT, Trader, Analyst} \iqfirm{any} \iqlevel{junior}]\label{iq:iv1:mental-technique:36}
Estimate $47 \times 62 \times 81 \times 19$ to within 1\%, without writing down any intermediate product. How do you correct for the rounding?
\end{question}

\begin{question}[$\star\star$ \iqroles{QR, QT, Trader, DS} \iqfirm{any} \iqlevel{junior}]\label{iq:iv1:mental-technique:37}
You toss a fair coin 20 times. Estimate the chance of exactly 10 heads in your head, without computing a binomial coefficient. Then estimate the chance of exactly 12 heads.
\end{question}

\begin{question}[$\star\star\star$ \iqroles{QR, QT} \iqfirm{market maker, hedge fund} \iqlevel{junior, experienced}]\label{iq:iv1:mental-technique:38}
How many ways are there to choose 25 items from 50? Give two significant figures without a calculator. Then do the same for choosing 30 from 100.
\end{question}

\begin{question}[$\star\star\star$ \iqroles{QR, QT, QD, Dev} \iqfirm{any} \iqlevel{junior, experienced}]\label{iq:iv1:mental-technique:39}
How many bits does it take to record the order of a shuffled 52-card deck, with the fewest bits possible? Estimate it without a calculator. How many decimal digits does $52!$ have?
\end{question}

\begin{question}[$\star\star\star$ \iqroles{QR, QT, Trader, Analyst, Risk} \iqfirm{bank, hedge fund} \iqlevel{junior, experienced}]\label{iq:iv1:mental-technique:40}
You receive \$1 at the end of each year for 30 years, and the discount rate is 7\% a year. Estimate the present value to one decimal place, without a calculator. Then price a 30-year bond that pays a 5\% coupon yearly, at the same rate.
\end{question}
